\begin{proof}
    First, it must be shown that $F( f ) \circ f \approx_E 1_n$ and $f \circ F( f ) \approx_E 1_m$ for each $f \in P( \Sigma )( n, m )$.
    This follows by structural induction on the morphisms in $P( \Sigma )$.
    \begin{itemize}
    \item \textbf{Identity}.
          Assume that $f = 1_n$.
          Since $F$ is a prop functor, then $F( f ) = F( 1_n ) = 1_n$.
          Then $F( f ) \circ f = 1_n \circ 1_n = 1_n$ and $f \circ F( f ) = 1_n \circ 1_n = 1_n$.
          Since $\approx_E$ is reflexive, then $F( f ) \circ f \approx_E 1_n$ and $f \circ F( f ) \approx_E 1_n$.
    \item \textbf{Symmetries}.
          Assume that $f = \sigma$.
          Since $F$ is a strict symmetric monoidal functor, then $F( \sigma ) = \sigma$.
          By definition, $F( f ) \circ f = \sigma \circ \sigma = 1_2$ and $f \circ F( f ) = \sigma \circ \sigma = 1_2$.
          Since $\approx_E$ is reflexive, then $F( f ) \circ f \approx_E 1_2$ and $f \circ F( f ) \approx_E 1_2$.
    \item \textbf{Generators}.
          Assume that $f = x$ for some $x \in \Sigma$.
          Then $n = \dom( x )$ and $m = \cod( x )$.
          Moreover, $F( f ) \circ f = F( x ) \circ x \approx_E 1_n$ and $f \circ F( f ) = x \circ F( x ) \approx_E 1_m$.
    \item \textbf{Sequential Composition}.
          Assume that $f = b \circ a$ for $a: n \to s$ and $b: s \to m$.
          Then by the inductive hypothesis, the following equations hold.
          \begin{align*}
                F( a ) \circ a &\approx_E 1_n
                &
                a \circ F( a ) &\approx_E 1_x
                &
                F( b ) \circ b &\approx_E 1_x
                &
                b \circ F( b ) &\approx_E 1_m
          \end{align*}
          Since $F$ is a contravariant functor, then the following equation holds.
          \begin{align*}
              F( f ) \circ f
                &=
                ( F( a ) \circ F( b ) ) \circ ( b \circ a )
                \approx_E
                F( a ) \circ 1_x \circ a
                =
                F( a ) \circ a
                \approx_E
                1_n
              \\
              f \circ F( f )
                &=
                ( b \circ a ) \circ ( F( a ) \circ F( b ) )
                \approx_E
                b \circ 1_x \circ F( b )
                =
                b \circ F( b )
                \approx_E
                1_m
          \end{align*}
          Then $F( f ) \circ f \approx_E 1_n$ and $f \circ F( f ) \approx_E 1_m$ by the transitivity of $\approx_E$.
    \item \textbf{Parallel Composition}.
          Assume that $f = a \boxtimes b$ for $a: s \to t$ and $b: x \to y$.
          Clearly $n = s + x$ and $m = t + y$.
          Then by the inductive hypothesis, the following equations hold.
          \begin{align*}
                F( a ) \circ a &\approx_E 1_n
                &
                a \circ F( a ) &\approx_E 1_x
                &
                F( b ) \circ b &\approx_E 1_x
                &
                b \circ F( b ) &\approx_E 1_m
          \end{align*}
          Since $F$ is a strict monoidal functor, then the following equation holds.
          \begin{align*}
              F( f ) \circ f
                &=
                ( F( a ) \boxtimes F( b ) ) \circ ( a \boxtimes b )
                =
                ( F( a ) \circ a ) \boxtimes ( F( b ) \circ b )
                \approx_E
                1_s \boxtimes ( F( b ) \circ b )
                \approx_E
                1_s \boxtimes 1_x
            \\
            f \circ F( f )
                &=
                ( a \boxtimes b ) \circ ( F( a ) \boxtimes F( b ) )
                =
                ( a \circ F( a ) ) \boxtimes ( b \circ F( b ) )
                \approx_E
                1_t \boxtimes ( b \circ F( b ) )
                \approx_E
                1_t \boxtimes 1_y
          \end{align*}
          Then $F( f ) \circ f \approx_E 1_s \boxtimes 1_x = 1_n$ and $f \circ F( f ) \approx_E 1_t \boxtimes 1_y = 1_m$ by the transitivity of $\approx_E$.
    \end{itemize}
    Then by the principle of structural induction, $F( f ) \circ f \approx_E 1_n$ and $f \circ F( f ) \approx_E 1_m$ for each $f \in P( \Sigma )( n, m )$.
    Next, let $( f, g ) \in E$ where $f: n \to m$.
    Then $f \approx_E g$, $F( f ) \circ f \approx_E 1_n$, and $g \circ F( g ) \approx_E 1_m$ by the previous result.
    Then the following equation holds.
    \begin{equation*}
        F( f )
        =
        F( f ) \circ 1_m
        \approx_E
        F( f ) \circ g \circ F( g )
        \approx_E
        F( f ) \circ f \circ F( g )
        \approx_E
        1_n \circ F( g )
        =
        F( g )
    \end{equation*}
    Then $F( f ) \approx_E F( g )$ by the transitivity of $( \approx_E )$.
    Then $\pi_E^\text{op}( F( f ) ) = \pi_E^{\text{op}}( F( g ) )$.
    Then by the universal property of prop quotients, there exists a unique prop functor $H: P( \Sigma ) / E \to ( P( \Sigma ) / E )^\text{op}$ such that $H \circ \pi_E = \pi_E^\text{op} \circ F$.
    It remains to be shown that $H$ is a dagger functor.
    To this end, let $f \in P( \Sigma )( n, m )$.
    Then $f \circ F( f ) \approx_E 1_m$.
    Likewise, $F( f ) \circ F^\text{op}( F( f ) ) \approx 1_n$.
    This means that $f = f \circ 1_n \approx_E f \circ F( f ) \circ F^\text{op}( F( f ) ) \approx_E 1_m \circ F^\text{op}( F( f ) ) F( F^\text{op}( f ) )$.
    Then it follows that $\pi_E( f ) = \pi_E( F^\text{op}( F( f ) ) ) = H^\text{op}( \pi_E^\text{op}( F( f ) ) ) = H^\text{op}( H( \pi_E( f ) ) )$.
    Since $f$ was an arbitrary, then $\pi_E = H^\text{op} \circ H \circ \pi_E$.
    Then by uniqueness, $H^\text{op} \circ H = 1_{P( \Sigma ) / E}$.
    Then $H$ is a dagger functor.
    Finally, let $f: n \to n$.
    Then $F( f ) \circ f \approx_E 1_n$ and $f \circ F( f ) \approx_E 1_m$.
    This means that $H( \pi_E( f ) ) \circ \pi_E( f ) = 1_n$ and $\pi_E( f ) \circ H( \pi_E( f ) ) = 1_n$.
    Since $f$ was arbitrary and $\pi_E$ is surjective, then $H( f ) \circ f = 1_n$ and $f \circ H( f ) = 1_m$ for each $f \in P( \Sigma ) / E$.
    In other words, every morphism in $P( \Sigma ) / E$ is unitary.
\end{proof}